\chapter{Kinematics and Stresses}

In this section, the kinematics and the different stress tensors used in this paper are defined. The relation between them is also given. Finally, the definition of the strain energy density is given as well as its relation with the different stress tensors. These definitions are used throughout the paper and are important for the understanding of the following sections.

\section{Kinematics}

The kinematics of a body is defined by the transformation between the initial configuration and the current configuration. This transformation can be defined as stated by Equation \ref{eq:transformation}. Thoughout this paper, the initial configuration will be denoted by capital letters while the current configuration will be denoted by lowercase letters.

\begin{equation} \label{eq:transformation}
    \x = \vec{\phi} (\X, \, t)
\end{equation}

Based on this transformation, it is possible to define the transformation gradient (Equation \ref{eq:F}) and the Jacobian of the transformation (Equation \ref{eq:J}).

\begin{equation} \label{eq:F}
    \F = \tensorTwoGradO{\x} = \pderiv{\x}{\X}
\end{equation}

\begin{equation} \label{eq:J}
    \J = \det{\F}
\end{equation}

Based on the tensor $\F$, it is possible to define multiple tensors to represent the defromation. In this codes, two of then are used: the right Cauchy-Green deformation tensor $\C$ and the Green-Lagrange strain tensor $\E$ as given by Equations \ref{eq:C} and \ref{eq:E} respectively.

\begin{equation} \label{eq:C}
    \C = \F^T \F
\end{equation}

\begin{equation} \label{eq:E}
    \E = \frac{1}{2} \left(\C - \tensorTwo{\mathrm{I}}\right) = \frac{1}{2} \left(\F^T \F - \tensorTwo{\mathrm{I}}\right)
\end{equation}

\section{Stresses}

Four different stress tensors are used throughout this paper and are listed below:

\begin{itemize}
    \item $\cauchy$: Cauchy stress tensor
    \item $\kirchhoff$: kirchhoff stress tensor
    \item $\pkOne$: First Piola-Kirchhoff stress tensor
    \item $\pkTwo$: Second Piola-kirchhoff stress tensor
\end{itemize}

They are all stress' representations but under different configurations. They can all be related under the following equation:

\begin{equation}
    \cauchy = \J^{-1} \cdot \kirchhoff = \J^{-1} \cdot \pkOne \F^T = \J^{-1} \cdot \F \pkTwo \F^T
\end{equation}

\section{Strain Energy Density}

The strain energy density is defined as:

\begin{equation}
    \potential (\F, \, \X) = \int\pkOne :\dot{\F} \dt
\end{equation}

From this definition, a new definition for $\pkOne$ can be stated:

\begin{equation}
    \pkOne = \pderiv{\potential}{\F}
\end{equation}

This can also be written as a function of the tensor $\C$ as given below:

\begin{equation}
    \potential (\F, \, \X) = \potential (\C, \, \X) = \int \frac{1}{2}\pkTwo : \dot{\C} \dt 
\end{equation}

This means the $\pkTwo$ can be written as:

\begin{equation}
    \pkTwo = 2 \pderiv{\Psi}{\C} = \pderiv{\Psi}{\E}
\end{equation}
